%% 296_Zeit_Offen_Invariant_En_ch.tex
%% FFGFT -- Dok. 296: Time, Openness and Invariance
%% Author: Johann Pascher
%% Date: July 2026
%% Language: English

\section*{Doc.~296 \quad Time, Openness and Invariance}

\subsection*{Observer, scale and the simplicity of the core}

Every observer sits on a scale. They do not see the object --- they see the
\emph{projection} of the object onto their scale. Projections from different
scales must differ, not because the observers err, but because the generative
object is richer than any single view. The perspective distortion is not a
mistake --- it is structural, the unavoidable consequence of the object
($T^4$, compact, boundaryless) not fitting into any single projection.
Hardly any scale can be surveyed simultaneously with another; the results of
observation therefore already differ perspectively as soon as a further
observer is added.

This is the deeper reason why quantum mechanics and the theory of relativity
are so difficult to reconcile: they are projections onto different scales of
the same object, and of course they look different --- not because one is
wrong, but because neither is complete.

All the more remarkable, therefore, is what follows: that the core is so
simple. One parameter $\xi = 4/30000$. One relation $\tilde{T} \cdot m = 1$.
One compact torus $T^4/\mathbb{Z}_3$. From these fall the fine-structure
constant $\alpha$, the lepton masses, the Koide scalar $Q = 2/3$, the phase
$\theta = 2/9$ --- the entire complexity of the Standard Model, the whole
richness of the phenomena that appear so overwhelming on the various observer
scales. The complexity does not reside in the core --- it resides in the
projections. The core is simple.

This is not self-evident. It could equally have been that nature starts afresh
on each scale, that there is no common root, that the variety is irreducible
--- that would be the generic outcome. That a single structure instead carries
all scales becomes visible only when one stops treating one's own scale as the
only possible one.

\subsection*{From body to society --- the same structure at every level}

That observers on different scales see different projections is not merely a
physical abstraction --- it begins with the simplest biological perception.
Two eyes are already two observers at slightly different positions. The
separation is small but real; the brain \emph{constructs} a three-dimensional
world from two slightly different projections, a world that exists nowhere
directly in that form. Depth is a reconstruction, not direct perception.
Further sense organs --- proprioception, balance, touch, hearing --- each
contribute their own scale and their own projection, and the synthesis of all
of them feels like immediate reality, even though it is a highly processed
achievement.

This has a direct consequence for physics: that we \emph{experience} the
world as three-dimensional is not evidence that it fundamentally is. It is the
finding that three dimensions are the projection onto which our sensory
apparatus has been optimised. $T^4$ beneath this does not contradict the
three-dimensional experience --- it is the reason why the projection is so
stable and consistent.

The same structure continues one level higher. A society is a collection of
observers, each on their own scale, each with their own projection. What
arises as social reality --- norms, values, worldviews, conflicts --- is the
synthesis of millions of slightly different projections, none of them
complete, none of them wrong in any absolute sense, but each perspectively
distorted. And as with binocular vision: the synthesis feels like direct
perception, even though it is a construction.

This explains several things that would otherwise be puzzling. Conflicts are
so persistent because two observers on different scales see the same object
and reach different conclusions --- both correct for their own projection,
both incomplete. As long as no common higher-level scale is accessible, the
conflict cannot be resolved by argument alone; one would have to shift the
projection plane, and that is the most difficult thing there is. Historical
change is therefore so slow: a society has calibrated its synthetic
achievement to a particular scale, which works well within that scale and
generates blindness to others. Paradigm shifts are scale shifts, and scale
shifts are painful because the previous coherent picture dissolves before the
new one has formed.

Science is the systematic attempt to triangulate projections --- multiple
observers, multiple scales, contradictions treated as signal rather than
noise. This is not self-evident; it runs against the natural impulse to take
one's own projection for reality.

The arc back to the core: that despite the multiplicity of projections ---
biological, physical, social --- something invariant underlies them all,
something that shows itself completely in no single projection but shines
through in all of them, is not self-evident. The generic outcome would be
that each scale had its own irreducible truth. That there is instead a common
core --- in physics $\xi$ on $T^4$, in experience the synthetic achievement
that carries all projections --- is the actual surprise, at every level.

\subsection*{Scale sequence as a succession of projection shifts}

From particle physics through atomic, crystalline and biological to galactic
structures, the same generative object appears fundamentally different at each
level --- not because the objects differ, but because each scale makes a
different symmetry of the core visible while integrating the others away.

At the \emph{particle scale} the projection cuts most finely: the
$\mathbb{Z}_3$ symmetry, the $\mathrm{SU}(3)$ colour structure and the
discrete masses are direct imprints of the $T^4$ geometry. The fine-structure
constant $\alpha$ and the phase $\theta = 2/9$ are most directly accessible
here --- the core is nearest on this scale. At the \emph{atomic scale}
($\sim a_0$) the colour structure has already been integrated away and appears
as neutral; what becomes visible are electromagnetic resonances --- orbitals
as standing waves, discrete energy levels as torus geometry on this coarser
scale. At the \emph{crystalline scale} ($\sim$\,Å--$\mu$m) the translation
symmetry of the lattice comes to the fore; collective modes (phonons, band
structures) arise as new effective degrees of freedom that do not exist at the
particle scale. At the \emph{biological scale} (nm--m) broken symmetries
dominate --- chirality, self-replication, information processing; life is a
coherence window, a scale at which a particular kind of complexity can just
remain stable, because below it thermal fluctuations and above it gravity
force other projections. At the \emph{galactic scale} (kpc--Mpc) everything
except gravity has been integrated away; the spiral structure of galaxies
carries the same logarithmic self-similarity as the K\textsubscript{frak}
recursion --- not coincidentally, but as an expression of the scale invariance
of the core at this projection level.

\subsection*{Scale transitions and their anomalies}

The transitions between scales are particularly instructive --- precisely
because anomalies appear there that cannot be explained from within a single
scale. A transition is where the previously dominant symmetry tips and a new
one comes to the fore; at this tipping point the system is simultaneously
shaped by both projections and fully described by neither. This structurally
generates anomalies: deviations that appear as disturbances from one side and
as harbingers of a new order from the other.

Familiar examples include phase transitions (critical fluctuations whose
correlation lengths diverge --- the system \enquote{senses} on all scales
simultaneously that a transition is imminent), the transition from classical to
quantum mechanics (decoherence as projection stabilisation, not as fundamental
collapse), or the transition from biological to social organisation (emergence
of collective intelligence not foreshadowed at the single-cell scale). In all
cases the anomaly lies \emph{in the transition}, not within either scale
itself --- and is therefore not fully explicable from any single scale.

For FFGFT this is a directed pointer: the places where a scale's standard
description reaches its limits --- unresolved anomalies, unexplained constants,
emergent phenomena without derivation --- are precisely the places where a
projection shift occurs. They are not errors of the theory but signatures of
the transition.

\subsection*{Speculation: unobservable intermediate ranges}

The scales considered so far are those accessible to us directly or indirectly
--- through instruments, through consequences, through indirect inference. But
the fact that scales are observer perspectives admits a further thought, which
is explicitly booked here as \emph{speculation} (P35: not a derived finding
but an open question):

Between the observable scales there may lie ranges that are in principle
inaccessible to us --- not for technical reasons but structurally, because
none of our sense organs or instruments projects onto these intermediate
scales. Projections are selective; they make one symmetry visible and suppress
others. It is not excluded that there are scales on which coherent structures
exist that leave no direct imprint in any of our accessible projections ---
at most an indirect one, as an anomaly at the edges of neighbouring scales.

This would not be mystical but structurally consistent: if the generative
object $T^4$ is richer than any projection, and if projections are always
selective, then the generic outcome is that there are regions of the object
captured by none of our projections. Dark matter and dark energy would, in
this reading, not be new substances but signatures of such intermediate
ranges --- structures at scales onto which our instruments do not project, but
which make themselves felt as anomalies at the edges of the galactic and
cosmic projections.

This remains speculation. But it is directed speculation: it names precisely
where to look for evidence --- at the transitions, in the anomalies, at the
places where a scale reaches its explanatory limits without internal reason.

\subsection*{Starting point: a recursion with no reachable boundary}

The K\textsubscript{frak} recursion
\begin{equation}
  \xi_{n+1} = \xi_n\,(1 - 100\,\xi_n)
\end{equation}
does not run endlessly in the sense of uniform repetition --- it
asymptotes. $\xi_n$ decays monotonically towards zero, but never reaches zero
in finitely many steps. This is not a gap in the construction; it is its
statement: the sequence is infinitely long, but the gaps between successive
steps shrink exponentially, precisely by the factor $1/e$ per $2\pi$ turn.

What does this mean for experienced, measurable time? The measurable time
step is not $n$ but $\rho = 1/(100\,\xi_n)$ --- and this reciprocal
\emph{grows}. The further the recursion proceeds, the larger $\rho$ becomes
and the coarser the time scale on which something changes. Experienced from
within, one does not encounter a uniform sequence but a \emph{slowing}
dynamics: early steps are dense and fast (small $\rho$, large $\xi$,
sub-Planck-proximate), late steps are wide and sluggish (large $\rho$, small
$\xi$, macroscopic). The experienced time is the integral over $\rho$ --- and
this integral grows logarithmically: $D(k) = \sum 100\,\xi_n \to
\ln(1 + k/75)$. The cumulative defect \emph{is} the measurable time evolution
(Doc.~295).

From this it follows immediately: experienced time has no end. $\rho$ grows
without bound, $D(k)$ diverges, $\xi_n$ never reaches zero. There is no final
step, no endpoint, no \enquote{afterwards}. And from the other side --- the
inner side --- the same holds in mirror image: $\xi$ never reaches its
starting value $\xi_0$ from below, $\rho$ never reaches $75 = 1/(100\,\xi_0)$
from above. No beginning, no end --- both sides are asymptotic limits, not
reachable boundaries.

\subsection*{Unrolled --- temporally continuous --- space curved ---
invariant}

These four properties are not accidentally connected; they follow from one
another.

\textbf{Unrolled.} The spiral unrolls the scale structure of $T^4$ into a
directed time coordinate. Not because time \enquote{really is} a spiral, but
because the projection --- the observer's position on the spiral --- singles
out a direction. The unrolling is the transition from the compact picture to
the open one: the same object, read from a particular scale.

\textbf{Temporally continuous.} This follows necessarily. The Hilbert-space
bijection $H = L^2(T^4) \otimes \mathbb{C}^3$ is lossless (Doc.~230/282,
bijection type~III, Doc.~270); what is discretely open remains continuously
open --- otherwise the translation would not be faithful but distorting. The
memory kernel (Fourier transform of the discrete spectral density
$\sum_k g_k^2\,\delta(\omega - \omega_k)$ of the $T^4$ connection-Laplacian,
Doc.~283) is defined on $\mathbb{R}$, not on an interval. No boundary on the
left, no boundary on the right. Revivals without end, BLP backflow without
closure --- the same open structure, now in the continuous picture.

\textbf{Space curved.} This is the $T^4$ core itself. The curvature is not
inserted after the fact; it is the generative structure: $\xi$ on
$T^4/\mathbb{Z}_3$ \emph{is} the curvature from which the lepton masses,
$\alpha$, and $\theta = 2/9$ follow (Doc.~006/011/291/293). The geometry is
not background --- it is the object. Compact means finite volume but
\emph{no boundaries}: a geodesic on the torus runs without end, never
encountering an edge.

\textbf{Invariant.} This is the deepest of the four points. The logarithmic
spiral does not change its size (Doc.~295): it is self-similar --- rotated by
one turn and scaled by $e$, it lies congruent with itself. Rotation and
scaling are \emph{the same} operation for it. \enquote{Inward by $1/e$} and
\enquote{outward by $e$} do not describe something growing or shrinking; they
only indicate which scale one is currently examining. The object is invariant
under this motion --- there is no absolute size, only ratios, and ratios are
the intrinsic quantities under $\tilde{T} \cdot m = 1$. What changes is the
\emph{observer scale}, not the object.

\subsection*{No beginning, no end --- and no coincidence problem}

The four properties together yield a consistent picture: an invariant, curved,
boundaryless object --- unrolled into an open, continuous time coordinate
through the act of projection. No beginning, no end, no absolute size, no
inserted curvature.

This differs structurally from models that require boundaries. A beginning at
$t = 0$ is not a prediction but a limit at which the description ceases to
hold --- the singularity is not a result but an admission. A finite time
interval requires a starting point, and the starting point demands a
justification: why now, why this state, why these initial conditions. The
coincidence problem, the flatness problem, the fine-tuning of initial
conditions --- all arise because a finite interval has a distinguished location.

FFGFT has no distinguished location. Every observer sits somewhere on an open
trajectory; \enquote{now} is the current scale, not a privileged moment. The
experienced time is not an interval between two points --- it is an open
trajectory on a boundaryless object.

\subsection*{New problems as a sign of progress}

Such a view generates new questions --- this is to be expected and is not an
objection. An approach that raises no new problems usually also explains
nothing new; it is without consequence. The decisive distinction is what
\emph{kind} of problems arise: whether they are \emph{directed follow-on
questions} (what still needs to be computed?) or genuine breaks (does
$\tilde{T} \cdot m = 1$ break?).

The exchange is asymmetric in FFGFT's favour: in place of unexplained
stipulations --- why these masses, why $\alpha = 1/137$, why $\theta = 2/9$,
where does non-locality come from --- there appear consequences of a
derivation, consequences that can be posed and whose resolution advances the
programme. New problems arising from a deep solution are the signature of
genuine progress, not of its absence. The task is not to avoid them but to
record clearly which are directed follow-on questions --- and so far, all of
them are.

\bigskip
\noindent\textit{Note~(P35):} The statements in this document are structural
in nature --- they follow from $\tilde{T} \cdot m = 1$, the recursion, and
the compactness of $T^4$. The cosmological transfer (open time, no
singularity) is a \emph{consequence} of this structure, not a standalone
cosmological model. As a consequence it must be free of contradiction with
the data --- not complete in the sense of covering all phenomena of a
dedicated model. This distinction between freedom from contradiction and
completeness is to be maintained.
